\section{Container Data Plane Integration}
\label{sec:data-plane}

The data plane runs the work that would make a control plane miserable: APIs, long-running services, data pipelines, optimization workloads, scheduled jobs, and integration workers. In the deployed system, that means EKS namespaces, Horizon/Grid-style module Lambdas, DataMap pipeline orchestration, Optima-style optimization services, Airflow DAGs, connector jobs, and Fargate tasks.

This is why the paper uses ``container data plane'' rather than ``Kubernetes backend.'' The execution fabric is mixed on purpose. Some work belongs in services, some in Lambdas, some in short-lived containers, some in scheduled pipelines, and some in connector workers. The governance rule should not change with the substrate.

\subsection{Representative execution surfaces}

Table~\ref{tab:execution_surfaces} summarizes the main surfaces governed by the control plane. The details differ; the contract does not. Each surface receives a small scope envelope, rejects missing scope, and avoids rediscovering tenant identity from user input.

\begin{table}[H]
  \centering
  \scriptsize
  \begin{tabularx}{0.94\textwidth}{p{0.18\textwidth} p{0.30\textwidth} X}
    \toprule
    \textbf{Surface} & \textbf{Workload shape} & \textbf{Control-plane responsibility} \\
    \midrule
    Horizon/Grid module Lambdas & Tenant-scoped query and analytics APIs & Gate module routes, inject \scopeprefix{}, preserve actor/effective context \\
    DataMap orchestration & Subject $\rightarrow$ Analysis $\rightarrow$ Graph pipelines and graph-version history & Attach tenant, user, subject, analysis, and version scope before pipeline execution \\
    Optima-style optimization & Long-running portfolio and scenario optimization jobs & Launch scoped compute and keep job intent tied to effective tenant context \\
    Airflow DAGs & Scheduled and ad-hoc data workflows & Receive scope and environment context rather than resolving membership inside DAG code \\
    Connector workers & External-system sync, webhook, and ingestion jobs & Bind connector metadata, credentials, and job events to tenant-owned records \\
    \bottomrule
  \end{tabularx}
  \caption{Representative data-plane surfaces. The substrate varies, but governance enters through the same control-plane boundary.}
  \label{tab:execution_surfaces}
\end{table}

\subsection{Module Lambdas: Horizon and Grid}

Horizon and Grid are useful examples because they look like ordinary product APIs until tenant scope gets involved. Horizon serves module-specific queries and derived views. Grid serves time-series analytics, predictions, and derived signals. Both sit behind module-gated routes in the control plane.

The control plane first resolves the actor and effective tenant, checks that the module is enabled, then delegates to the module Lambda with an explicit scope envelope. The downstream Lambda treats \scopeprefix{} as required input. If it is absent, the request fails before any domain logic runs.

\begin{lstlisting}[style=python, caption={Simplified module Lambda scope check.}]
class TenantScopedHandler:
    def __init__(self, query_params):
        self.schema_prefix = query_params.get("schema_prefix")
        if not self.schema_prefix:
            raise BadRequestError("schema_prefix is required")
        validate_schema_prefix(self.schema_prefix)
        self.core_schema = f"{self.schema_prefix}_core"
\end{lstlisting}

The Lambda does not decide whether \code{acme} is the right tenant. It only checks that the scope is present and syntactically safe. The authority for deriving \code{acme} remains in the control plane. That division is subtle but important. If every module Lambda is allowed to rediscover tenant membership, the boundary is gone.

\subsection{Scope envelope}

The envelope passed downstream is intentionally small. It carries just enough information for execution to be correct and auditable, not enough to turn every service into a policy interpreter.

\begin{table}[H]
  \centering
  \scriptsize
  \begin{tabularx}{0.94\textwidth}{p{0.22\textwidth} p{0.26\textwidth} X}
    \toprule
    \textbf{Field} & \textbf{Typical consumers} & \textbf{Purpose} \\
    \midrule
    \code{actor_id} & Audit logs, support workflows & Preserves the real authenticated identity \\
    \code{effective_tenant} & All tenant-scoped services & Names the tenant context in effect after optional \viewas{} resolution \\
    \code{effective_role} & Module APIs, UI contracts & Carries the role used for capability checks \\
    \scopeprefix{} & Horizon/Grid, query services, jobs & Selects the tenant database namespace, e.g., \texttt{\{prefix\}\_core} \\
    \code{modules} & API filtering, route guards & Declares enabled product modules for the effective context \\
    \code{subject_id}, \code{analysis_id}, \code{version_id} & DataMap & Locates a stable artifact hierarchy for graph refinement and history \\
    \code{connector_id}, \code{job_id} & Connector workers & Binds external-system sync to tenant-owned metadata and credentials \\
    \bottomrule
  \end{tabularx}
  \caption{The scope envelope. The control plane derives it; data-plane systems consume it.}
  \label{tab:scope_envelope}
\end{table}

The point is not to standardize these exact field names forever. The point is to stop hiding tenancy in query parameters, environment variables, README files, and tribal memory.

\subsection{DataMap as a scoped artifact pipeline}

DataMap has a different shape from Horizon and Grid. It is closer to an interactive modeling studio: a user starts from a subject, creates analyses, generates or refines causal graphs, and keeps history across versions. The workflow is naturally iterative. A graph is not overwritten casually; it is refined, versioned, and attached to the subject and analysis that produced it.

Figure~\ref{fig:datamap_flow} shows the artifact hierarchy. The control plane does not need to understand graph semantics. It does need to attach the right tenant and user context before the pipeline creates, reads, or updates artifacts.

\begin{figure}[H]
  \centering
  \resizebox{0.86\textwidth}{!}{%
  \begin{tikzpicture}[
    node distance=12mm and 18mm,
    box/.style={draw, rounded corners=3pt, align=center, minimum width=28mm, minimum height=11mm, font=\small\sffamily},
    arrow/.style={-{Latex[length=2.3mm]}, thick, color=gray!75},
    note/.style={font=\scriptsize\sffamily, align=center}
  ]
    \node[box, fill=blue!10, draw=blue!55!black] (tenant) {Tenant scope\\\scopeprefix{}};
    \node[box, fill=green!10, draw=green!55!black, right=of tenant] (subject) {Subject\\business object};
    \node[box, fill=green!10, draw=green!55!black, right=of subject] (analysis) {Analysis\\question/context};
    \node[box, fill=orange!15, draw=orange!60!black, right=of analysis] (graph) {Graph\\current version};
    \node[box, fill=violet!12, draw=violet!55!black, below=16mm of graph] (history) {History\\prior versions};
    \node[box, fill=gray!10, draw=gray!60!black, below=16mm of analysis] (source) {Data source\\dataset now, connector later};

    \draw[arrow] (tenant) -- (subject);
    \draw[arrow] (subject) -- (analysis);
    \draw[arrow] (analysis) -- (graph);
    \draw[arrow] (source) -- (analysis);
    \draw[arrow] (graph.south west) .. controls +(-8mm,-4mm) and +(-8mm,4mm) .. node[note,left] {save\\version} (history.north west);
    \draw[arrow] (history.north east) .. controls +(8mm,4mm) and +(8mm,-4mm) .. node[note,right] {restore} (graph.south east);
  \end{tikzpicture}%
  }
  \caption{DataMap artifact flow. Tenant scope is attached before the pipeline touches subjects, analyses, graph versions, or connector-backed inputs.}
  \label{fig:datamap_flow}
\end{figure}

A typical storage layout mirrors the same hierarchy. The exact table and bucket names are deployment details, but the key pattern is stable.

\begin{lstlisting}[style=python, caption={Representative DataMap artifact key pattern.}]
subject_key  = f"tenant#{tenant_id}#subject#{subject_id}"
analysis_key = f"{subject_key}#analysis#{analysis_id}"
graph_key    = f"{analysis_key}#graph#{graph_id}#version#{version}"
object_key = (
    f"tenants/{tenant_id}/subjects/{subject_id}/"
    f"analyses/{analysis_id}/graphs/{graph_id}/{version}.json"
)
\end{lstlisting}

That shape is deliberately boring. It makes version history auditable, lets multiple analyses exist under one subject, and leaves room for data sources to evolve from uploaded datasets to connector-backed queries without changing the tenant boundary.

\subsection{Connector ingestion}

Connector ingestion adds one more wrinkle: external systems. A connector has tenant-owned metadata, tenant-owned credentials, webhooks or polling jobs, and downstream events. The control plane owns the connector record and launches or receives work under tenant scope. Connector workers execute sync logic; they do not decide which tenant owns a token or destination.

\begin{figure}[H]
  \centering
  \resizebox{0.88\textwidth}{!}{%
  \begin{tikzpicture}[
    node distance=12mm and 18mm,
    box/.style={draw, rounded corners=3pt, align=center, minimum width=31mm, minimum height=11mm, font=\small\sffamily},
    arrow/.style={-{Latex[length=2.3mm]}, thick, color=gray!75}
  ]
    \node[box, fill=yellow!18, draw=yellow!55!black] (admin) {Tenant admin\\console action};
    \node[box, fill=blue!10, draw=blue!55!black, right=of admin] (cp) {Control plane\\connector API};
    \node[box, fill=violet!12, draw=violet!55!black, above right=of cp] (meta) {DynamoDB\\metadata + status};
    \node[box, fill=red!8, draw=red!55!black, below right=of cp] (secrets) {Secrets Manager\\tokens};
    \node[box, fill=green!10, draw=green!55!black, right=34mm of cp] (worker) {Connector worker\\sync / webhook};
    \node[box, fill=orange!12, draw=orange!55!black, right=of worker] (events) {Tenant-scoped\\jobs + events};

    \draw[arrow] (admin) -- (cp);
    \draw[arrow] (cp) -- (meta);
    \draw[arrow] (cp) -- (secrets);
    \draw[arrow] (cp) -- node[above, font=\scriptsize] {launch scoped job} (worker);
    \draw[arrow] (worker) -- (events);
    \draw[arrow] (worker) |- (meta);
  \end{tikzpicture}%
  }
  \caption{Connector ingestion boundary. The control plane owns tenant metadata and credential binding; workers perform sync and event handling inside the injected tenant scope.}
  \label{fig:connector_boundary}
\end{figure}

This is the same pattern again. Interactive APIs, modeling pipelines, optimization jobs, and connector workers all receive scope. None of them should need a private theory of tenancy.

\subsection{Database scoping}

Database credentials are managed by AWS Secrets Manager~\cite{aws_secrets_manager} and reached through RDS Proxy~\cite{aws_rds_proxy}. Per-tenant isolation uses PostgreSQL schemas: module Lambdas set \code{search_path} to a tenant-specific schema such as \code{acme_core} before executing queries~\cite{postgresql_set}.

PostgreSQL does not allow parameter binding for \code{SET search_path}, so \scopeprefix{} is validated before use. The control plane validates before forwarding. Module Lambdas re-validate before issuing the statement. This second check is defense in depth. It protects against malformed strings. It does not replace authoritative derivation.

\subsection{On-demand jobs}

Fargate jobs receive scope at launch time through the ECS \code{RunTask} payload~\cite{ecs_runtask}. The job does not perform a new user lookup and does not resolve tenant membership. It receives a scoped job intent and runs.

That rule keeps batch work inside the same governance model as interactive APIs. Jobs are not a back door. They are data-plane executions with a larger time budget.

\subsection{One-directional dependency}

The direction of dependency is deliberate. The control plane calls into the data plane. The data plane does not call back to ask who it is allowed to be. This avoids a circular governance path where every service is both execution engine and policy interpreter.

The resulting contract is simple enough for engineers to remember: if a service needs tenant scope, it receives it from the control plane. If a service needs to invent tenant scope, the design has already failed.
